\section{Global obstructions for conditioned approximate factors}
\label{app:approximation-refinements}
The local rank estimates of Section~\ref{sec:conditioned-approximation}
have global consequences when the factors have fixed labels and vary
smoothly over the entire contact sphere. These hypotheses are not supplied
by an arbitrary approximate lift. We first retain tangent-bundle information
without requiring the selected factors to lie on cone boundaries. We then
obtain stronger channelwise conclusions for nonzero boundary selections.

\subsection{A robust tangent splitting}
Let $n\geq1$ and suppose proper cones $K_i$ of dimensions $m_i\geq2$
admit globally $C^2$ maps $A_i:\Sph^n\to K_i$,
$B_i:\Sph^n\to K_i^*$ satisfying
\begin{equation}\label{eq:global-approximate-slack}
 \sum_i\ip{A_i(x)}{B_i(y)}=h_C(y)-\ip xy,\qquad
 \B_2^{n+1}\subset C\subset(1+\epsilon)\B_2^{n+1}.
\end{equation}
Put $r_i=m_i-2$ and assume $\sum_i r_i=n$. At $x$ write
$a_i=A_i(x)$, $b_i=B_i(x)$, $X_i=dA_i(x)$, $Y_i=dB_i(x)$.
Assume uniformly that
$\norm{a_i},\norm{b_i}\geq\mu>0$ and
$\norm{X_i},\norm{Y_i}\leq L$. Fix $0<r_0<\pi$ and $H>0$ such that,
for every unit $u\in T_x\Sph^n$ and $|t|\leq r_0$,
\[
 |\ip{(A_i\circ\exp_x(tu))''}{b_i(x)}|\leq H,\qquad
 |\ip{a_i(x)}{(B_i\circ\exp_x(tu))''}|\leq H.
\]
The opposite factor is held fixed at the center in these inequalities.

\begin{theorem}[Approximate saturation gives a bundle splitting]
\label{thm:robust-splitting}
If $\epsilon<\mu^2$, $\sqrt{2\epsilon/H}\leq r_0$, and
$\mathcal E_k(\epsilon)<1$, where $k$ is the number of blocks and
$\mathcal E_k$ is defined in~\eqref{eq:uniform-smooth-error}, then
\[
 T^*\Sph^n=\bigoplus_{i:r_i>0}E_i,\qquad \rank E_i=r_i.
\]
Consequently, the number of $m_i=3$ blocks is at most $\rho(n+1)-1$.
If all blocks have dimension three, $n\in\{1,3,7\}$.
If $n\geq2$ is even, exactly one block has positive $r_i$, and
that block has dimension $n+2$.
\end{theorem}
\begin{proof}
Let $\delta_i=\ip{a_i}{b_i}$. Equation~\eqref{eq:global-approximate-slack}
gives $\delta_i\geq0$, $\sum_i\delta_i\leq\epsilon$, and
$-\sum_iX_i^*Y_i=I$. Since
$\gamma_i=\delta_i/(\norm{a_i}\norm{b_i})<1$, the vectors $a_i,b_i$
are linearly independent. The spaces
$W_i(x)=a_i(x)^\perp\cap b_i(x)^\perp$ therefore form continuous
bundles of rank $r_i$. Define the canonical bilinear forms
\[
 R_i(u,v)=-\ip{P_{W_i}X_i u}{P_{W_i}Y_i v}.
\]
They depend continuously on $x$ and have rank at most $r_i$.
On $W_i$, both $P_{b_i^\perp}P_{a_i^\perp}$ and $P_{W_i}$ are the
identity. On the remaining plane their difference has norm $\gamma_i$.
Thus the proof of Lemma~\ref{lem:near-complementarity} gives
$\norm{-X_i^*Y_i-R_i}\leq e_i$, with the sharpened
error~\eqref{eq:smooth-factor-error}. It follows that
$\norm{I-\sum_iR_i}\leq\mathcal E_k(\epsilon)<1$.
The sum is invertible and
\[
 n\leq\sum_i\rank R_i\leq\sum_i r_i=n.
\]
All ranks attain their bounds, and their images are in direct sum.
Continuity and constant rank give the stated subbundles.
For $n\geq2$, every real line bundle on $\Sph^n$ is trivial, so
rank-one summands give independent vector fields; Adams's bound is
$\rho(n+1)-1$~\cite{Adams1962}. For $n=1$ the tangent bundle is
already trivial. A complete rank-one splitting forces the familiar
parallelizable dimensions $1,3,7$. For even $n$, a splitting into two
positive intermediate ranks would give
$e(T\Sph^n)=e(E)e(F)=0$, since intermediate integral cohomology vanishes.
This contradicts Euler number two.
\end{proof}

The stronger subset-rank restrictions in
Proposition~\ref{prop:abstract-splitting-ranks} also apply. The canonical
maps in this proof are essential: unrelated pointwise low-rank
approximants do not automatically define vector bundles. The hypotheses
exclude rays with nonzero selected factors, since for a ray
$\delta_i=\norm{a_i}\norm{b_i}\geq\mu^2$. Allowing a factor to vanish
or losing uniform control is a permitted failure of the hypotheses.

\subsection{An invariant of fixed block reparameterizations}
There is a coordinate-independent version of this sufficient obstruction,
although it is still an invariant of the chosen selections and contact
parameterization. Retain the global slack identity and the saturated
budget $\sum_i(m_i-2)_+=n$, now permitting rays. For a fixed family
$S_i\in GL(m_i)$, transform
\[
 A_i\mapsto S_iA_i,\qquad B_i\mapsto S_i^{-T}B_i,
 \qquad K_i\mapsto S_iK_i.
\]
At each center compute $U_i^S,V_i^S,\alpha_i^S,\beta_i^S,\gamma_i^S$
as in Lemma~\ref{lem:near-complementarity} for the transformed maps,
using geodesics of a fixed radius $r_0$. In particular,
\[
 \alpha_i^S=\frac{\psi_{r_0}(\delta_i,H_{A,i})}{\norm{S_i^{-T}b_i}},
 \qquad
 \gamma_i^S=\frac{\delta_i}{\norm{S_i a_i}\norm{S_i^{-T}b_i}}.
\]
The contracted derivative bounds $H_{A,i},H_{B,i}$ are unchanged by
this transformation. Let $e_i^S$ be~\eqref{eq:smooth-factor-error},
or $\alpha_i^S\beta_i^S$ for a ray, and set
\begin{equation}\label{eq:gauge-degeneration-index}
 \mathfrak D=\inf_{S\in\prod_i GL(m_i)}
 \max\left\{\sup_{x,i}\gamma_i^S(x),\,
                 \sup_x\sum_i e_i^S(x)\right\}.
\end{equation}
If a base factor vanishes anywhere, the displayed value for that gauge
is defined as $+\infty$.

\begin{corollary}[Gauge-invariant degeneration]
\label{cor:gauge-degeneration}
If $T^*\Sph^n$ admits no splitting with ranks $(m_i-2)_+$, then
$\mathfrak D\geq1$.
\end{corollary}
\begin{proof}
If $\mathfrak D<1$, some fixed gauge has both suprema below one.
Its $\gamma_i<1$ excludes rays and makes the primal and dual factors
nonparallel everywhere. The canonical maps of the preceding proof
then have rank at most $m_i-2$ and an invertible sum, supplying the
forbidden splitting. If the original coordinates are changed by $T_i$,
all subsequent gauges $S_i$ range through the same net transformations
$S_iT_i$ as before. Thus the infimum is invariant under every fixed
blockwise invertible reparameterization. No covariance of the individual
Euclidean projectors is needed.
\end{proof}
This index does not assert a lower bound on a barrier Hessian or a
Newton-system condition number. The infimum ranges over fixed linear
block gauges, not parameter-dependent gauges or changes of the contact
sphere itself.

\subsection{Rank-one phase dropout}
The next statement is stronger than the frame consequence in source
sphere dimensions three and seven, but requires nonzero boundary maps.
Let $M$ be a closed, connected, simply connected smooth Riemannian manifold.
Let $A,B:M\to\partial Q_3\setminus\{0\}$ be $C^1$ and write
\[
 A=a(1,p),\qquad B=b(1,-q),\qquad a,b>0,\quad p,q:M\to\Sph^1.
\]
Set $\sigma=\ip AB$, $C_x(u,v)=-\ip{dA_xu}{dB_xv}$ and
\[
 K_A=\sup_M\norm{d\log a},\qquad K_B=\sup_M\norm{d\log b},
 \qquad K=\min\{K_A,K_B\}.
\]
Bilinear forms need not be symmetric; all norms are operator norms.

\begin{theorem}[Quantitative phase dropout]\label{thm:robust-phase}
If $0\leq\sigma\leq\epsilon$ and $\norm{d\sigma}\leq\eta$, then
there is $x_*\in M$ with
\[
 \norm{C_{x_*}}\leq K\eta+K^2\epsilon.
\]
If every $C_x$ is within $\delta$ of a positive rank-one form whose
nonzero eigenvalue is at least $\mu>0$, then
\begin{equation}\label{eq:phase-threshold}
 \delta+K\eta+K^2\epsilon\geq\mu.
\end{equation}
\end{theorem}
\begin{proof}
Simple connectivity lifts $p$ to a real $C^1$ phase. At a maximum of
that phase, $dp=0$, so $dA=\alpha\otimes A$, where
$\alpha=d\log a$. At this point
\[
 \ip A{dB}=d\sigma-\sigma\alpha,
 \qquad C=-\alpha\otimes(d\sigma-\sigma\alpha).
\]
Its norm is at most $K_A\eta+K_A^2\epsilon$. A critical point of the
dual phase gives $C=-(d\sigma-\sigma\beta)\otimes\beta$, with
$\beta=d\log b$, and the bound using $K_B$. Taking the smaller bound
proves the first assertion. At the corresponding point, any rank-one
approximant of norm at least $\mu$ has distance at least
$\mu-\norm C$ from $C$, proving~\eqref{eq:phase-threshold}.
\end{proof}

The hypotheses can be checked using norms of the original factors:
$\norm A=\sqrt2a$ and
$\norm{dA(u)}^2=2|da(u)|^2+a^2\norm{dp(u)}^2$, whence
$K_A\leq L_A/m_A$ if $\norm A\geq m_A$ and $\norm{dA}\leq L_A$.
The dual bound is identical. Reciprocal constant scalar gauges leave
$K_A,K_B$, $\sigma$, and $C$ unchanged.

\begin{corollary}[Saturated approximate phase channels]
\label{cor:saturated-phase}
Let $M=\Sph^n$, $n\geq2$, and let $n$ pairs satisfy the preceding
theorem with constants $\epsilon_i,\eta_i,K_i$. Suppose
\[
 \left\|\sum_iC_i-I\right\|\leq\xi,\qquad
 \norm{C_i-R_i}\leq\delta_i,
\]
where $R_i$ is positive of rank one at every point; the choices of
$R_i$ need not be continuous. Put $\Delta=\sum_i\delta_i$ and
$E_i=K_i\eta_i+K_i^2\epsilon_i$. Then, for every $i$,
\[
 \xi+\Delta+\delta_i+E_i\geq1.
\]
In particular, common bounds give $\xi+(n+1)\delta+E\geq1$.
\end{corollary}
\begin{proof}
There is nothing to prove if $\xi+\Delta\geq1$. Otherwise
$\norm{\sum_iR_i-I}\leq\xi+\Delta<1$. Write
$R_i=\lambda_i\omega_i\otimes\omega_i$ with $\norm{\omega_i}=1$.
At any point, choose a unit vector killed by the at most $n-1$
covectors $\omega_j$, $j\ne i$. Evaluation on that vector gives
$\lambda_i\geq1-\xi-\Delta$. Apply Theorem~\ref{thm:robust-phase}
to the $i$th channel with this spectral floor.
\end{proof}

We record explicitly how first-derivative moduli supply $\eta$.
For $r_0$ below the injectivity radius, suppose covectors $d\sigma$
at endpoints of geodesic segments of length $t\leq r_0$ differ by
at most $\omega(t)$ after parallel transport. Then
\begin{equation}\label{eq:modulus-interpolation}
 \norm{d\sigma}\leq
 \inf_{0<s\leq r_0}
 \left\{\frac{\epsilon}{s}+
              \frac1s\int_0^s\omega(t)\,dt\right\}.
\end{equation}
Indeed, along a unit geodesic in the gradient direction, the integral
of the derivative is at least $s\norm{d\sigma}-\int_0^s\omega(t)dt$,
whereas the change in $\sigma$ is at most $\epsilon$.
If $\omega(t)\leq Ht$, this is $\psi_{r_0}(\epsilon,H)$.
For $H>0$ with $\sqrt{2\epsilon/H}\leq r_0$, it gives
$\eta\leq\sqrt{2H\epsilon}$. If $H=0$, $d\sigma$ is locally parallel;
its vanishing at a maximum and connectedness give $d\sigma=0$
everywhere. A common modulus tending to zero at zero makes the
right side of~\eqref{eq:modulus-interpolation} tend to zero with
$\epsilon$.

In a nonnegative kernel $F(x,y)=\sum_i\ip{A_i(x)}{B_i(y)}$ with
$F(x,x)\leq\epsilon$, every individual defect satisfies
$0\leq\sigma_i\leq\epsilon$. Thus the preceding inequalities apply
separately to each channel when it has its own derivative modulus.
A bound on the derivative of the sum does not replace these individual
bounds: derivative cancellation is possible.

\subsection{Higher-rank dropout from forbidden sphere submersions}
Let $1\leq r<n$, let $K\subset E$ be proper of dimension $r+2$,
and choose $\ell\in\intr K^*$. Suppose the normalized base
$D=\{z\in K:\ell(z)=1\}$ has a $C^1$ boundary $J$.
Let $A:\Sph^n\to\partial K\setminus\{0\}$ and
$B:\Sph^n\to\partial K^*\setminus\{0\}$ be $C^1$. Define
\[
 \begin{gathered}
 a=\ell(A),\quad p=A/a:\Sph^n\to J,\quad
 K_A=\sup\norm{d\log a},\\
 g_B(v)=\ip A{dB(v)},\quad \eta_B=\sup\norm{g_B},
 \quad C=-dA^*dB.
 \end{gathered}
\]
Write $s_r(C)$ for the $r$th largest singular value.

\begin{theorem}[Robust submersion obstruction]
\label{thm:robust-submersion}
If $(n,r)\notin\{(3,2),(7,4),(15,8)\}$, there is a point with
\[
 s_r(C)\leq K_A\eta_B.
\]
Consequently, if every $C$ is within $\delta$ of a rank-$r$ form
whose $r$th singular value is at least $\mu>0$, then
\begin{equation}\label{eq:submersion-threshold}
 \delta+K_A\eta_B\geq\mu.
\end{equation}
In a Hopf dimension pair, the same conclusion holds if the particular
map $p$ has a critical point. If a dual normalized base also has $C^1$
boundary, the term $K_A\eta_B$ in~\eqref{eq:submersion-threshold}
may be replaced by $\min\{K_A\eta_B,K_B\eta_A\}$, with dual definitions.
\end{theorem}
\begin{proof}
Radial projection from a relative interior point of the convex base
identifies $J$ $C^1$ diffeomorphically with $\Sph^r$. The inverse is
$C^1$ because the supporting normal at each boundary point has strictly
positive inner product with its radial vector. Outside the three Hopf
pairs, Proposition~\ref{prop:sphere-product-submersions} rules out an
everywhere submersive map $p$. Hence $\rank dp\leq r-1$ somewhere.
Writing $\alpha=d\log a$, differentiation of $A=ap$ gives
\[
 C(u,v)=-a\ip{dp(u)}{dB(v)}-\alpha(u)g_B(v).
\]
At that point the first term has rank at most $r-1$ and the second
has norm at most $K_A\eta_B$. The singular-value min--max inequality
gives $s_r(C)\leq K_A\eta_B$. Singular values change by at most the
operator perturbation norm, proving~\eqref{eq:submersion-threshold}.
Applying the same argument to the normalized dual map supplies the
second bound, possibly at a different point; both inequalities hold,
so their minimum is valid.
\end{proof}

Here $g_B(x)$ is the one-sided derivative
$d_y\ip{A(x)}{B(y)}|_{y=x}$, with $x$ held fixed. It can be bounded
from small contact defect and a modulus in that variable. Specifically,
if $F(x,y)=\ip{A(x)}{B(y)}\geq0$, $F(x,x)\leq\epsilon$, and the
second-variable derivative has uniform geodesic modulus $\omega_B$,
then
\begin{equation}\label{eq:one-sided-interpolation}
 \eta_B\leq\inf_{0<s\leq r_0}
 \left\{\frac{\epsilon}{s}+
              \frac1s\int_0^s\omega_B(t)\,dt\right\}.
\end{equation}
Move the second argument from $x$ in the direction opposite to
$g_B(x)$; the endpoint is nonnegative, so integration gives this
inequality. The first-variable bound for $\eta_A$ is identical.
Lipschitz constants $H_A,H_B$ give the corresponding $\psi$ bounds,
and, when their optimizing radii fit,
\[
 \delta+\min\{K_A\sqrt{2H_B\epsilon},
                    K_B\sqrt{2H_A\epsilon}\}\geq\mu.
\]
The finite-radius formula~\eqref{eq:one-sided-interpolation} remains
valid when either Lipschitz constant is zero.

A total diagonal derivative bound is insufficient for this higher-rank
argument. With $\sigma=\ip AB$,
\[
 g_B=d\sigma-\sigma\,d\log a-a\,dp^*B.
\]
At a rank-one phase critical point $dp=0$, recovering
Theorem~\ref{thm:robust-phase}. For $r>1$ a critical point only forces
$\rank dp<r$ and the last term need not vanish. If it is bounded by
$\tau$ and $\norm{d\sigma}\leq\eta$, the valid replacement is
\[
 \delta+K_A\eta+K_A^2\epsilon+K_A\tau\geq\mu.
\]

Finally, consider a saturated profile $\sum_i r_i=n$ with at least
two positive ranks, where every block satisfies these boundary and
base-regularity hypotheses and its channel is uniformly $\delta_i$
close to rank $r_i$ with singular-value floor $\mu_i$. At least one
pair $(n,r_i)$ is forbidden: outside source dimensions $3,7,15$ all
proper pairs are forbidden; in those three dimensions the only
permitted proper targets are $2,4,8$, respectively, and copies of that
target dimension cannot sum to $n$. Thus, under common bounds
$\delta_i\leq\delta$, $K_{A,i}\leq K$, $\eta_{B,i}\leq\eta_\partial$,
and $\mu_i\geq\mu$, every such profile satisfies
\[
 \delta+K\eta_\partial\geq\mu.
\]
A single full-capacity block is outside this obstruction.

These conclusions combine elementary perturbation estimates with the
classical sphere topology already cited in
Appendix~\ref{app:topology-refinements}. The added statements concern
quantitative degeneration of specified global cone-contact selections.
They do not follow from the finite slack-matrix approximation framework
of~\cite{GPTApprox2015}, and do not obstruct arbitrary nonsmooth,
unlabelled, or vanishing-factor approximations.
